\SGASubsection{7.5}{}
\SGARefTarget{sga3:xxvi:subsec:7:5}

Let \(P,P_1\) be two minimal parabolic subgroups of \(G\), with
Levi subgroups \(L,L_1\), and let \(Q,Q_1\) be the largest central
split subtori of \(L,L_1\), respectively.  Then the pairs
\((P,Q)\) and \((P_1,Q_1)\) are conjugate: there is a \(g\in G(S)\)
such that
\[
 \Int(g)P=P_1,\qquad \Int(g)Q=Q_1.
\]
Indeed, \(P\) and \(P_1\) are conjugate by \SGARefLink{sga3:xxvi:corollary:5-7}{5.7}, so we may suppose
\(P=P_1\).  Then \(L\) and \(L_1\) are conjugate by an element of
\(P(S)\), by \SGARefLink{sga3:xxvi:cor:1:8}{1.8}.  Moreover, if \(g,g'\) are two elements of \(G(S)\)
which conjugate \((P,Q)\) to \((P_1,Q_1)\), then \(g'^{-1}g\)
normalizes \(P\) and \(Q\), hence also \(P\) and \(L\).  But
\[
 \Norm_G(P)\cap\Norm_G(L)
 =P\cap\Norm_G(L)
 =L
 =\Centr_G(Q).
\]
Consequently the isomorphism
\[
 Q\xrightarrow{\ \sim\ }Q_1
\]
induced by \(\Int(g)\) is independent of \(g\).

Let \(\mathcal R,\mathcal R_1\) be the root data defined by \SGARefLink{sga3:xxvi:theorem:7-4}{7.4} in
\[
 \operatorname{Hom}_{S\text{-}\mathrm{gr}}(Q,\mathbf G_{m,S})
 \quad\text{and}\quad
 \operatorname{Hom}_{S\text{-}\mathrm{gr}}(Q_1,\mathbf G_{m,S}),
\]
and let \(R^+,R_1^+\) be the systems of positive roots corresponding
to \(P,P_1\).  The canonical isomorphism \(Q\simeq Q_1\) just
constructed carries \((\mathcal R,R^+)\) to
\((\mathcal R_1,R_1^+)\).  We may therefore define the
\emph{based relative root datum} of \(G\) over \(S\), by identifying
the various pairs \((\mathcal R,R^+)\) through the transitive system
of isomorphisms above.

From now on, denote this based root datum by
\[
 \mathcal R(G/S)
 =(\underline M,\underline M^*,\underline R,
   \underline R^*,\underline R^+).
\]
For every pair \(P\supset Q\) as above there is therefore a
canonical isomorphism
\[
 \underline M\xrightarrow{\ \sim\ }
 \operatorname{Hom}_{S\text{-}\mathrm{gr}}(Q,\mathbf G_{m,S})
\]
which carries \(\underline R\), respectively \(\underline R^+\), to
the roots of \(G\), respectively of \(P\), relative to \(Q\), and
carries \(W(\mathcal R)\) to \(W_G(Q)(S)\).

\SGASubsection{7.6}{}
\SGARefTarget{sga3:xxvi:subsec:7:6}

Let again \(Q\) be a maximal split torus of \(G\), let \(P\) be a
minimal parabolic subgroup of \(G\) with Levi subgroup
\(\Centr_G(Q)\), let
\[
 (M,M^*,R,R^*,R^+)
\]
be the corresponding based root datum of \SGARefLink{sga3:xxvi:theorem:7-4}{7.4}, and let \(\Delta\)
be the set of simple roots in \(R^+\).  For every
\(A\subset\Delta\), let
\[
 R_A=R^+\cup(\mathbf Z A\cap R^-).
\]
Thus \(R_A\) consists of the positive roots together with the
negative roots which are integral linear combinations of elements
of \(A\).  It is a closed set of roots, by \SGARefLink{sga3:XXI:3.1.4}{Exposé XXI, 3.1.4}, and
every closed set containing \(R^+\) is uniquely of this form, by
\SGARefLink{sga3:XXI:3.3.10}{Exposé XXI, 3.3.10}.

By \SGARefLink{sga3:xxvi:proposition:6-1}{6.1} there is a unique smooth subgroup \(P_A\) of \(G\), with
connected fibers, which contains \(\Centr_G(Q)\) and satisfies
\[
 \operatorname{Lie}(P_A)
 =\operatorname{Lie}(G)^0
  \oplus\bigoplus_{\alpha\in R_A}\operatorname{Lie}(G)^\alpha.
\]
The following result now follows immediately from \SGARefLink{sga3:xxvi:proposition:6-1}{6.1}, the
conjugacy of minimal parabolic subgroups, and the fact that the set
of roots of a parabolic subgroup of \(G\) containing \(Q\) is
closed \textup{(which follows from \SGARefLink{sga3:xxvi:prop:1:3}{1.4} after splitting)}.

\begin{sgastatement}{\SGAInlineRefTarget{sga3:xxvi:proposition:7-7}Proposition 7.7}
\begin{enumerate}[label=\textup{(\roman*)}]
\item The map
\[
 A\longmapsto P_A
\]
is an inclusion-preserving bijection from the set of subsets of
\(\Delta\) onto the set of parabolic subgroups of \(G\) containing
\(P\).
\item Every parabolic subgroup of \(G\) is conjugate by an element
of \(G(S)\) to a unique \(P_A\).
\end{enumerate}
\end{sgastatement}
\SGARefTarget{sga3:xxvi:prop:7:7}

\SGASubsection{7.8}{}
\SGARefTarget{sga3:xxvi:subsec:7:8}

Let \(P\supset Q\) be as above.  Consider the relative root datum of
\(G\) over \(S\) defined in \SGARefLink{sga3:xxvi:subsec:7:5}{7.5} and the canonical isomorphism
\[
 f:(M,M^*,R,R^*,R^+)
   \xrightarrow{\ \sim\ }
   (\underline M,\underline M^*,\underline R,
    \underline R^*,\underline R^+).
\]
The set \(\Delta\) of simple roots in \(R^+\) is carried to the set
\(\underline\Delta\) of simple roots in \(\underline R^+\).
Thus every subset \(A\subset\Delta\) is carried to a subset
\[
 f(A)\subset\underline\Delta.
\]

\begin{sgastatement}{Definition 7.9.0}
Let \(H\) be any parabolic subgroup of \(G\).  By
\SGARefLink{sga3:xxvi:proposition:7-7}{7.7}\textup{(ii)}, it is conjugate to a unique \(P_A\).  Put
\[
 t_r(H)=f(A)\subset\underline\Delta.
\]
The conjugacy theorems show immediately that \(t_r(H)\) is
independent of the choice of the pair \(P\supset Q\).  It is called
the \emph{relative type} of \(H\).
\end{sgastatement}
\SGARefTarget{sga3:xxvi:def:7:9:0}

\begin{sgastatement}{\SGAInlineRefTarget{sga3:xxvi:proposition:7-9}Proposition 7.9}
\begin{enumerate}[label=\textup{(\roman*)}]
\item The map \(H\mapsto t_r(H)\) induces a bijection from the set of
\(G(S)\)-conjugacy classes of parabolic subgroups of \(G\) onto the
set of subsets of \(\underline\Delta\).
\item Let \(H\) be a parabolic subgroup of \(G\), \(P\) a minimal
parabolic subgroup contained in \(H\), \(Q\) the largest central
split torus of a Levi subgroup of \(P\), \(\Delta\) the set of
simple roots of \(P\) relative to \(Q\), and
\[
 f:\Delta\xrightarrow{\ \sim\ }\underline\Delta
\]
the canonical isomorphism.  For every \(\alpha\in\Delta\),
\[
 f(\alpha)\in t_r(H)
 \quad\Longleftrightarrow\quad
 \operatorname{Lie}(H)^{-\alpha}\ne0,
\]
and
\[
 H=P_A,\qquad A=f^{-1}\bigl(t_r(H)\bigr).
\]
\item If \(H,H'\) are two parabolic subgroups of \(G\) containing
\(P\), then \textup{(see \SGARefLink{sga3:xxvi:rem:3:8:1}{3.8.1}\textup{(ii)} and
\SGARefLink{sga3:xxvi:corollary:5-5}{5.5}\textup{(i)} for other equivalent conditions)}
\[
 t_r(H)\subset t_r(H')
 \quad\Longleftrightarrow\quad
 t(H)\subset t(H').
\]
\end{enumerate}
\end{sgastatement}
\SGARefTarget{sga3:xxvi:prop:7:9}
